\documentclass[10pt]{article}
\usepackage[margin=20mm]{geometry}
\usepackage{amsmath,amssymb,fontspec,unicode-math,xurl,hyperref}
\IfFontExistsTF{DejaVu Serif}{\setmainfont{DejaVu Serif}}{\setmainfont{Latin Modern Roman}}
\setmathfont{latinmodern-math.otf}
\setlength{\parindent}{0pt}\setlength{\parskip}{6pt}\sloppy
\begin{document}
Abstract: The proceedings version’s core filter-containment derivations are invalid: a finite analytic bounded filter cannot equal a nonconstant polynomial on an interval, and the proof’s single-sinusoid replacement fails on an admissible three-node graph. Two release-level evaluation defects are also established. Approximation with increasing capacity and permutation equivariance survive the declared checks. Scientific adjudication remains pending.\par
Author: AI-assisted technical review\par
Date:\par
Title: Technical Verification of GrokFormer\par
\# Identity and scope\par
The reviewed paper is Guoguo Ai, Guansong Pang, Hezhe Qiao, Yuan Gao and Hui Yan, GrokFormer: Graph Fourier Kolmogorov-Arnold Transformers, ICML 2025, PMLR 267, I25-43630 (local 0261) [Retained latex: \textbackslash{}cite\{paper\}]. The proceedings PDF has 20 pages, including appendices. All page numbers below are one-based physical PDF pages and agree with the printed pages. The released source is pinned to commit `014944357beaec6ab8dc45951466c56b734588d8` [Retained latex: \textbackslash{}cite\{code\}]; it is not attested as the experiment-time revision.\par
We checked the filter construction, Propositions 4.1–4.4 and their Appendix C proofs, experimental protocol, and relevant released source. One Major issue affects the core filter-containment derivations; two Minor issues affect the released evaluator. Approximation with increasing capacity and scalar spectral permutation equivariance survive the checks. No empirical ranking was reproduced or refuted.\par
\# Definitions and surviving arguments\par
Equation reference: eq:h. For the paper’s undirected binary graph with positive degrees, \[L=I-D^{-1/2}AD^{-1/2}=U\Lambda U^\top,\qquad 0\leq\lambda_i\leq2.\] For finite positive integers \(K,M\) and finite real coefficients, Eq. (9) defines \[
h(\lambda)=\sum_{k=1}^{K}\alpha_k\sum_{m=0}^{M}
\{a_{km}\cos(m\lambda^k)+b_{km}\sin(m\lambda^k)\}.\] The convolution is \(Uh(\Lambda)U^\top X\). Exact inclusion on \([0,2]\), approximation as \(M\) grows, and interpolation on one finite graph are different assertions.\par
For \(K=1\), the real trigonometric polynomials contain constants, form an algebra, and separate points on \([0,2]\): distinct points in an interval shorter than \(2\pi\) cannot share both sine and cosine. Their uniform closure is therefore \(C([0,2])\). Thus Proposition 4.4’s approximation conclusion can hold when \(M\) depends on the target and tolerance, even though its proof unnecessarily assumes additional differentiability. No fixed-\(M\) guarantee follows. Also, \[h(PLP^\top)=Ph(L)P^\top\] for a permutation \(P\). Equal eigenvalues receive equal responses, so eigenbasis signs and rotations within repeated eigenspaces cancel. This verifies the stated scalar convolution, not every truncated implementation.\par
\# Substantive findings\par
Problem 1 (Invalid core filter-containment derivations: Major).  Propositions 4.2–4.3 and their proofs do not establish the special-case relationships invoked to support the central expressiveness comparison.\par
\#\# Original claim and dependency\par
Page 5 states that polynomial and Specformer filters are simplified variants of the proposed filter, using these special cases to support greater universality and flexibility. The theoretical comparison is an explicit contribution on p. 2. Appendix C, pp. 18–19, Eqs. (13) and (15)–(23), supplies the derivations. This is a core theoretical dependency, irrespective of whether benchmark scores survive.\par
\#\# Exact polynomial inclusion\par
The proof reaches the polynomial Eq. (16) by removing the nonlinear functions. That changes the architecture instead of identifying parameters of Eq. [eq:h] (\#eq:h). Every finite instance of Eq. [eq:h] (\#eq:h) is real analytic and satisfies \[|h(t)|\leq\sum_{k=1}^{K}|\alpha_k|\sum_{m=0}^{M}(|a_{km}|+|b_{km}|),\quad t\in\mathbb R.\] It cannot equal the admissible nonconstant polynomial \(p(\lambda)=\lambda\) on \([0,2]\): equality on an open interval would extend to all real arguments by analytic uniqueness, contradicting boundedness. The continuation argument tests agreement inside the graph-spectrum interval; it does not assume that graph eigenvalues lie outside it.\par
\#\# A finite, fully admissible check of Eq. (19)\par
The proof of Proposition 4.3 uses \(a_0\lambda=R\sin(\lambda+\phi)\) in Eqs. (19)–(20). Set \(a_0=1\) and take the three-node path: \[A=\begin{pmatrix}0&1&0\\1&0&1\\0&1&0\end{pmatrix},
\quad D=\operatorname{diag}(1,2,1),\quad
L=\begin{pmatrix}1&-1/\sqrt2&0\\-1/\sqrt2&1&-1/\sqrt2\\0&-1/\sqrt2&1\end{pmatrix}.\] The graph is undirected, binary, and has no zero degree. Its characteristic polynomial is \((1-t)((1-t)^2-1)\), giving normalized-Laplacian eigenvalues \(\{0,1,2\}\). All graph assumptions used by the paper are satisfied.\par
Write the proposed replacement as \(g(\lambda)=a\sin\lambda+b\cos\lambda\). Matching \(0\) gives \(b=0\); matching \(1\) gives \(a=1/\sin1\). At \(2\) the replacement yields \[g(2)=\frac{\sin2}{\sin1}=2\cos1=1.08060461173628\neq2.\] The contradiction therefore holds on a legal finite spectrum, not only on a continuous interval.\par
The text also calls this replacement approximate. Even under that interpretation, a single shifted sinusoid cannot attain arbitrary precision: every such \(g\) satisfies \(g(2)-2\cos1\,g(1)+g(0)=0\). Setting \(e_j=g(j)-j\) gives \[\max_{j=0,1,2}|e_j|\geq\frac{2-2\cos1}{2+2\cos1}
=\tan^2(1/2)=0.298446410409525.\] This bound diagnoses Eq. (19) only. It does not constrain the full multi-harmonic GrokFormer.\par
\#\# Counterevidence and bounded conclusion\par
Approximation with increasing \(M\) and interpolation on a particular finite spectrum remain possible. Interpreting a simplified variant as replacement of basis functions by the identity describes a modified architecture, not the claimed exact inclusion. In addition, the Specformer proof restricts attention to the identity and compares logarithmically spaced real frequencies \(\epsilon/10000^{i/M}\) with integer frequencies without a mapping. The restricted calculation cannot establish the unrestricted family’s inclusion.\par
The demonstrated invalid core derivation supports Major. It does not establish that the method is ineffective or experimentally inferior. The density statement in Section 2 is a supported replacement approximation claim, but it does not repair the original comparative containment theorem. Evidence: original-page rendering `s19.png`, independent `check.py`/`check.json`, and `theory\_notes.json`.\par
Problem 2 (Penn94 runs reuse the first split: Minor).  Appendix A.2.2, p. 14, relates five runs to five official Penn94 splits. The released `run.sh:7` requests five runs, and `train.py`, lines 129–135, varies seeds, but `utils.py:43–46` always reads split-array element `[0]`.\par
For five distinct saved partitions \(S_0,\ldots,S_4\), the five calls all select \(S_0\). Initialization may vary, but the literal split index is seed-independent. This does not implement an average over the five official partitions. If the stored partitions were identical, partition diversity would still not be established.\par
This is a release-level discrepancy. The appendix’s explicit five-split exception, rather than the main text’s generic ten-split wording, is the appropriate comparison. Historical code and per-run split logs are unavailable; the published Penn94 values are not thereby declared wrong. Passing and recording the intended split index would address the selection defect, but no author-code patch or new experiment was executed. Evidence: `code\_static.json` and the pinned source lines.\par
Problem 3 (Zero metrics when the epoch budget expires: Minor).  In `train.py:51–52`, the returned metrics are initialized to zero; only the patience-triggered branch at lines 88–92 assigns them. Line 95 returns them after either termination mode.\par
Use the allowed 2000-epoch budget from Appendix A.2.2. A finite control-flow witness has validation loss \(\ell_t=1/(t+1)\), \(t=1,\ldots,2000\), validation accuracy \(0.5\), and test accuracy \(0.8\). Every loss improves on the initial threshold \(100\) and its predecessor, so the patience counter stays zero. The loop exhausts its budget and returns zero metrics. Either validation-selected checkpoint has test accuracy \(0.8\).\par
This independently constructed trace checks the reporting branch; it does not assert that a particular optimization run produced that trajectory. If every published run early-stopped, this defect would not affect those numbers. Thus the supported impact is Minor for the release, with historical impact unresolved. Moving selection after both termination modes repairs this reporting failure for nonempty completed runs; a complete training repair was not verified.\par
\# Other checked scope and remaining work\par
The inspected source implements the stated trigonometric family and spectral multiplication order. Its node-level attention uses an unnormalized bilinear product, whereas Eq. (2) defines softmax attention. These operations are not identical, but the text explicitly introduces an efficient-attention variant and separately specifies vanilla attention for graph classification. This supports a scope clarification, not another Major finding.\par
Tables 2–6 and the experimental appendix were inspected for claim correspondence. Training, data validation, benchmark regeneration and historical experiment provenance were not checked empirically. Most original figure panels were inspected only through captions and extracted descriptions, not quantitatively audited visually. Later revisions, external errata and the full Specformer paper were not independently surveyed. Separate OpenReview attachments remain unenumerated after the acquisition record’s HTTP 403 response. None of these gaps is used as evidence of an error.\par
\# Conclusion and state\par
Dimension | Recorded state\par
Supported assessment | Major core derivation; two Minor release defects\par
Scientific adjudication | Pending for all findings\par
Repair verification | Approximation statement justified; containment not restored; no code patch run\par
Author adoption | Not requested; unknown\par
Remaining work | Historical split/run provenance, empirical effect, revised containment argument\par
9 G. Ai, G. Pang, H. Qiao, Y. Gao and H. Yan. GrokFormer: Graph Fourier Kolmogorov-Arnold Transformers. PMLR 267 (ICML 2025), proceedings PDF, 20 pages. https://proceedings.mlr.press/v267/ai25d.html. OpenReview: `wFBBh8bcoC`. GGA23/GrokFormer. Public repository snapshot, commit `014944357beaec6ab8dc45951466c56b734588d8`, dated 2025-07-10. https://github.com/GGA23/GrokFormer. Static inspection only.\par
\# Evidence, reproducibility and roles\par
Original PDF SHA-256: 3148572b588fa159d43f873e8c5251a44038e594a3718b2e7639909feb19f72b. Evidence names are relative to `verification/0261/`; original PDF and manifest reside in `papers/0261/`, while the pinned archive, member map and hierarchy reside in `sources/0261/`. `source\_check.json` preserves manifest, CRC, Git-blob and extracted-byte validation. Page text uses explicit physical-page ranges with controls removed only in derivatives.\par
The parent AI independently integrated the claims and calculations, with native auxiliary agents for source integrity, theoretical cross-checking and static-code checking. The parent compiled the exact final TeX with existing pdfLaTeX and inspected all resulting PDF pages; corresponding commands, logs, hashes and visual-check records accompany the report. The exposed system label is GPT-6 Astra Pro; the requested label is 6 Pro. An independently observed interface label and conversation URL are unavailable and remain null in provenance. No human adjudication, author contact, author-code execution, installation or empirical reproduction occurred. This is the initial review, not separate batch acceptance.\par
\end{document}
